%!TEX root = ../template.tex
%%%%%%%%%%%%%%%%%%%%%%%%%%%%%%%%%%%%%%%%%%%%%%%%%%%%%%%%%%%%%%%%%%%%
%% abstract-pt.tex
%% NOVA thesis document file
%%
%% Abstract in Portuguese
%%%%%%%%%%%%%%%%%%%%%%%%%%%%%%%%%%%%%%%%%%%%%%%%%%%%%%%%%%%%%%%%%%%%

\typeout{NT FILE abstract-pt.tex}%



A modelação de software exige uma compreensão aprofundada do domínio da solução, bem como das especificações de requisitos. Os modelos de design resultantes servem como “blueprints” para a fase de implementação, definindo a arquitetura do sistema, os requisitos de hardware e software, e a estrutura global do sistema. Para automatizar a criação destes modelos de design, foram propostas várias abordagens ao longo dos anos, incluindo estratégias baseadas em desenvolvimento orientado por modelos, regras e padrões linguísticos, sistemas baseados em regras e técnicas estatísticas ou de aprendizagem automática. Mais recentemente, os Large Language Models (LLMs) ganharam destaque. Embora estes modelos apresentem frequentemente resultados promissores, carecem muitas vezes de consciência contextual suficiente — um fator essencial para transformar corretamente as especificações de requisitos em modelos, como por exemplo diagramas de classes. Esta dissertação irá recorrer a técnicas de prompt engineering e a uma framework de Retrieval-Augmented Generation para reforçar a consciência contextual e a precisão dos LLMs na geração de modelos a partir de uma especificação de requisitos. Como primeiro passo, foi realizado um estudo de mapeamento sistemático para compreender melhor o estado atual do conhecimento. Com base nessas conclusões, começámos a experimentar várias técnicas de prompt engineering. O nosso objetivo é desenvolver um método aprimorado que aprofunde a compreensão contextual dos LLMs e facilite a criação automática de modelos de design a partir de especificações de requisitos.

% E agora vamos fazer um teste com uma quebra de linha no hífen a ver se a \LaTeX\ duplica o hífen na linha seguinte se usarmos \verb+"-+… em vez de \verb+-+.
%
% zzzz zzz zzzz zzz zzzz zzz zzzz zzz zzzz zzz zzzz zzz zzzz zzz zzzz zzz zzzz comentar"-lhe zzz zzzz zzz zzzz
%
% Sim!  Funciona! :)

% Palavras-chave do resumo em Português
% \begin{keywords}
% Palavra-chave 1, Palavra-chave 2, Palavra-chave 3, Palavra-chave 4
% \end{keywords}
\keywords{
  LLM \and
  Modelação de requisitos \and
  Retrieval-Augmented Generation\and
  diagramas UML\and 
  GAI
}

% to add an extra black line
